\section{A companion equation of degree 190 in ten positive integers}
\label{sec:positive190}

The compact companion has ten strictly positive integer variables, exact
total degree $190$, and $119$ live binary arithmetic operations
($74$ multiplications, $24$ additions, and $21$ subtractions).
Its coefficient tuple is unchanged from source 172.  This is a
positive-domain result, separate from the unrestricted-integer principal
construction.  It lowers the inspected ten-positive-variable degrees
$214$, $208$, and $206$, at a higher arithmetic cost than the $87$- and
$90$-operation alternatives.  It does not establish global minimality.

The new step is a transport polynomial $L=UEb^2$ that combines the
free-quotient chart with the whole-product center.  Transport by $bV$
would leave a degree-$197$ cross term.  Bare $b^2$ does not have the needed
uniform bound when the main Pell index is three.  The additional factor
$UE$ handles this exceptional index while keeping the cross term below
degree $190$.

\subsection{Complete formula}

The independent variables are
\[
 J,F,\alpha,i,T,w,s,n,\eta,\zeta>0.
\]
Use the complete inherited fixed tuple $(\beta,K_0,M_C,M_F)$, with
$\beta=2^{590069}-1$.  Its four full numerical payloads and finite
reconstruction recipe accompany the equation.  Define
\begin{align*}
 h&=\beta J,&q&=h+1,&C&=q-F-\alpha,\\
 X&=qw,&Y&=s,&Y'&=qY,\\
 k&=\eta+\zeta,&c&=kY+\eta,&A&=Y'(X+1)+1,\\
 \Delta&=A^2-1,&H&=4A-5,&v&=X(Y')^2,\\
 \ell&=vk,&E&=1+\ell(\ell+k),&U&=ic^2,\\
 a&=\Delta U^2E,&F_0&=1+a,&G&=EF_0,\\
 f&=1+2a,&Q&=4\Delta aF_0,&K&=X(3c-2X)+2.
\end{align*}
Also set
\begin{align*}
 R&=[h(q-F)+\alpha](q^2-1)+(M_C+qM_F)J,\\
 B&=(K_0+w)C-F,\qquad b=cT-R,\qquad V=fb-c,\\
 e_{\rm clr}&=6(k-R-1-q^4Y)-\ell(k^2-4),\\
 D&=6q^7,\qquad M=hDH.
\end{align*}
The center is the following division-free polynomial:
\begin{align*}
 e_2&=(kY')^2,\\
 S&=(kY')[\ell+k-(kY')(Y'+1)],\\
 \theta&=U^2b,\qquad R_c=\theta(3E+e_2)-(b-c),\\
 S_c&=4E^2+3S^2+e_2^2,\\
 \Omega&=U\{2SR_c+A(4EV-\theta S_c)\}.
\end{align*}
Finally put
\begin{align*}
 L&=UEb^2,\\
 Z&=hD\{H(n-LB-\Omega)+K\}-H(DB+h e_{\rm clr}),\\
 \mathcal P_{190}&=Q(Z^2-16M^2GV^2)+(16M^2G-Z^2).
\end{align*}
The equation is $\mathcal P_{190}=0$.  All named registers are
polynomial abbreviations, not extra quantified variables.  There is no
variable exponent, quotient, factorial, or hidden predicate in this
definition.

\subsection{Uniform transport and soundness}

The unchanged strong-root and all-index folding argument gives, at a zero,
positive integers $g,y$ with
\[
 G=g^2,\quad QV^2-(Q-1)y^2=1,\quad Z/M=\pm4gy,
 \quad c=\psi_A(p),\quad p\ge3\text{ odd}.
\]
These conclusions precede every sign assertion concerning $B,C,R,b,V$.
Indeed $Q=4\Delta^2U^2G$ and
$(Q-1)Z^2=16M^2G(QV^2-1)$ force the square class of $G$ by prime
valuations; the fixed square factor $16$ changes no valuation parity.
The retained adjacent-coefficient descent proves the initially rational
$y=|Z/M|/(4g)$ integral.  The signed folding then gives
$c\ge c_3=4\Delta+3$.

The case $b=0$ would give
$0<y^2-c^2=(c^2-1)/(Q-1)<1$, which is impossible.  Thus $b\ne0$,
$V$ has its sign, and the auxiliary odd-index classification gives
$|V|\ge4Q-3$.  Since $2c<f$ and $c_3\le c$,
\[
 |b|\ge\frac{|V|-c}{f}
 \ge4\Delta f-\frac{c+c_3}{f}
 >4\Delta f-1>3\Delta f.
\]
This explicitly permits auxiliary index three.  Also
\[
 y<\tfrac32|V|,\qquad |V|<\tfrac32f|b|,
 \qquad (AUE)^2-G=E(U^2E-1)>0,
\]
so $g<AUE$.

For the signed center estimate, temporarily regard $A$ as an indeterminate.
Then
\begin{align*}
 E&=e_2A^2+e_1A+e_0,\\
 e_1&=k^2Y'(1-2Y'-2(Y')^2),\\
 e_0&=1+k^2Y'(Y'+1)[Y'(Y'+1)-1].
\end{align*}
Let $Q_j$ be the formal polynomial quotient of $EV$ by $A^j$.  Exact
polynomial algebra gives
\[
 \Omega=U(4AEV-2Q_1-Q_3/2).
\]
Every coefficient of $Q_3$ is even, because
$EV=(b-c)E+2U^2b(A^2-1)E^2$ and $E$ has $A$-degree two.
The formula above is proof notation for the displayed division-free center.

Put $\delta=(Y'+1)/A<1/16$.  Absolute-coefficient evaluation gives
$E^\#(A)<2E$ and
\[
 (EV)^\#(A)
 <2E(|b|+c)+8(A^2+1)U^2E^2|b|<10|EV|.
\]
For the last inequality use $c+1<2U^2E$ and
$|V|>(f-c)|b|>\Delta U^2E|b|$; the resulting ratio is at most
$8+20/\Delta<10$.  Hence
\[
 |\Omega-4AUEV|
 <10U|EV|\left(\frac2A+\frac1{2A^3}\right)<AUE|V|.
\]
In particular $\Omega$ has the sign of $V$ and
$|\Omega|<5AUE|V|$.  Combining these estimates,
\[
 4gy+|\Omega|
 <\frac{33}{2}AUEf|b|
 <\frac{11A}{2\Delta}L<\frac L4.
\]
The final inequality follows from $A>q^2\ge1024$, and is independent
of the signs still to be recovered.

For proof only define
\[
 \rho=\frac{k-R-1-vk(k^2-4)/6}{q^7},\quad
 \sigma=\frac{Y}{q^3},\quad\kappa=\frac KH,
 \quad t=\frac Bh,\quad W=1+hL.
\]
The literal carrier satisfies
\[
 Z/M=n+\kappa+\sigma-Wt-\rho-\Omega.
\]
Since $k\ge2$ and $k-R-1<b$, one has $\rho<|b|/q^7$ for every
integer $R$.  If $B\le-1$, then $Wt<-L$, and therefore
\[
 n=\pm4gy+\Omega-\kappa-\sigma+Wt+\rho
 <L/4-L+|b|/q^7<0,
\]
contrary to $n>0$.  This proves $B\ge0$, hence $C>0$.

The authenticated mask now gives $3q<R<q^4-2$.  The unchanged normalized
index and signed auxiliary-index lemmas identify $p=R$.  Because $\Omega$
is an integer polynomial, $Z/M+\Omega$ is integral; because $W=1+hL$,
its residue modulo $h$ is unchanged.  The source-172 inverse Lucas argument
thus gives $m=(R+1)/2$, followed by the strict Pell projection $X=2^R$.
Only then is $q$ known to be a power of two, licensing the full
$q^7$ divisibility of the inverse correction.  The carrier finally yields
$q^3\mid Y$ and $h\mid B$.  The retained native reconstruction and the
whole coherent tuple's refutation calibration finish soundness.  In
particular no factor of three is canceled before the dyadic conclusion.

\subsection{Completeness and exact degree}

A finite refutation supplies the source-172 canonical positive outer data,
with $b,V>0$, integral $t\ge0$, integral $\rho,\kappa,\sigma$, and
$g>\kappa+\sigma+|\rho|$.  The center is positive on these data.
Choose
\[
 n=4gy+\Omega+(1+hL)t+\rho-\kappa-\sigma>0.
\]
Then $Z/M=4gy$, so $\mathcal P_{190}=0$.  This constructs fresh positive
witnesses; it does not assume that an arbitrary integer shear preserves
positivity.  Both directions are ordinary arguments formalizable in ZFC,
using the explicitly retained arithmetic and refutation calibration.

Let $Z_0$ be the carrier with $\Omega$ omitted and
$D_{\rm norm}=\Omega^2-16GV^2$.  The exact decomposition is
\[
 \mathcal P_{190}
 =Q[Z_0^2-2MZ_0\Omega+M^2D_{\rm norm}]
   +16M^2G-(Z_0-M\Omega)^2.
\]
In $\mathbb Z[A,Y',k,U,b,c]$ with weights $(4,2,1,5,4,2)$, the center
has $84$ monomials and weight $59$; the norm defect has $594$ monomials
and weight $94$, with complete highest part
\[
 [D_{\rm norm}]_{94}
 =8A^6(Y')^8k^8U^6b^2[A^2-6(Y')^4].
\]
The actual degrees of $(Q,M,G,V,Z_0,\Omega,L)$ are
$(72,12,46,36,41,59,27)$.  The five summands in the decomposition have
bounds $154,184,190,70,142$, respectively.  Put $q_0=\beta J$ and
$k_0=\eta+\zeta$.  The complete actual highest form is
\[
 [\mathcal P_{190}]_{190}
 =9\cdot2^{11}q_0^{74}w^{18}s^{52}i^{10}k_0^{32}
   (q_0-F)^2(w^2-6s^2).
\]
Its coefficient at $J^{74}w^{20}s^{52}i^{10}\eta^{32}F^2$ is
$9\cdot2^{11}\beta^{74}\ne0$.  Thus the degree is exactly $190$ for
the actual fixed tuple.  The independent checker collects all center and
defect coefficients, verifies the graph transformation and the leading
form, reconstructs all four fixed numerals, and proves genuine input
occurrence by nonzero modular derivatives.  These finite checks supplement
the unbounded proof and do not decide consistency.

For an explicit unrestricted-integer conversion of this companion, replace
each positive variable $x_j$ by $1+a_j^2+b_j^2+c_j(c_j+1)$.  The
three-square theorem makes this a surjection from $\mathbb Z^3$ onto the
positive integers: represent $4m+1$ by three squares and separate its
unique odd coordinate.  The resulting polynomial has $30$ unrestricted
variables and degree $380$.  This corollary is not an unrestricted-arity
record.
